\section*{Chapter \ref{ch:hf:a-venue-playbook} --- A Venue Playbook}

\begin{solution}{exo:hf:a-venue-playbook:1}
$0.125\times1{,}200+5=155$ orders. The runs with a pacing of eight seconds or less had 151 to 154 accepted: the budget, all of it, whatever was sent.
\end{solution}

\begin{solution}{exo:hf:a-venue-playbook:2}
Two orders every $2/0.125=16$ seconds. The home venue's quoting has no pacing at all (a gap of zero), which the profile's validation refuses.
\end{solution}

\begin{solution}{exo:hf:a-venue-playbook:3}
The top order gets $\min(0.4\times1{,}000,\,200)=200$. The other 800 contracts are shared among the remaining 4,800 in proportion to size: the
quoter gets $\lfloor 500\times800/4{,}800\rfloor=\lfloor83.3\rfloor=83$, and the rounding remainders go FIFO.
\end{solution}

\begin{solution}{exo:hf:a-venue-playbook:4}
Every message of every agent is delayed by the same 350 microseconds, so the arbitrageur's order still arrives before the quoter's cancel, and the
tape's flow, whose events are tens of milliseconds apart, meets the quoter's orders in the same order as without the delay. An asymmetric delay
that spares cancels, or a quoter that reacts to a signal the arbitrageur does not have, would change the races.
\end{solution}

\begin{solution}{exo:hf:a-venue-playbook:5}
The delay protects resting orders from the arbitrageur but applies to the quoter's own new orders: 52\% of its volume there was taken, not made,
because an order priced for the book it saw landed five seconds later in a book that had moved and crossed it. Those takes are at stale prices, in
the direction the price had already moved.
\end{solution}

\begin{solution}{exo:hf:a-venue-playbook:6}
In a uniform-price batch the quoter's bid at the touch executes only when the batch clears at or below it, which happens when the batch's selling
outweighs its buying: its fills come when the price is going its way least. It is never picked off because its jump cancel lands in the same interval
as the arbitrageur's order and is processed before the clearing. The fix is in the pricing, not the parameters: quote around a fair value that
anticipates the batch's clearing price, and size by the imbalance.
\end{solution}

\begin{solution}{exo:hf:a-venue-playbook:7}
An eight-second gap: 10,100 units against 8,333 at sixteen seconds, with 28 orders rejected in twenty minutes. The sixteen-second rule leaves
tokens unused (118 of about 155 accepted); the eight-second pacing uses them all, relying on the burst to absorb the moments when both sides
need a requote at once, and pays for it with a few rejects.
\end{solution}

\begin{solution}{exo:hf:a-venue-playbook:8}
Validation checks the profile against its own rules, and the simulator checks the strategy against the simulator's model of the venue: neither
checks the profile against the venue. The first days must run in the venue's test environment and then at small size under risk limits, comparing
fills, rejects and fees with the profile, before full size.
\end{solution}

\begin{solution}{pb:hf:a-venue-playbook:1}
\begin{enumerate}
\item See \cref{def:hf:a-venue-playbook:profile}; the venue half (what the venue does) and the quoting half (what the desk does there).
\item Matching (\cref{ch:hf:futures-market-making}), tick and lot, fees and unit (\cref{ch:hf:rebates-inverted-venues-and-tiers}), feed
(\cref{ch:hf:the-business-of-market-making}), order types (\cref{ch:hf:the-quoting-engine}), delays (\cref{ch:hf:latency-arbitrage-and-its-defence}),
budgets (\cref{ch:hf:crypto-high-frequency-trading-on-centralised-venues}), session and quirks.
\item Order-by-order feeds with new references on replace; a symmetric 350-microsecond delay and a disapproved asymmetric one; top order then pro
rata on some futures; message-efficiency ratios; execution-based risk tools on options; 100 orders per 10 seconds and 200,000 a day on a crypto
venue; a one-to-twelve-second in-play delay on placing bets.
\item So that the list of venues is complete and a strategy that assumes an order book is refused where there is none.
\item Size a multiple of the lot; post-only only where offered; post-only orders where a delay spares them; requotes paced within the order budget;
free cancels under a budget.
\item The crypto budget (no pacing) and the FX lot (100 units against a million).
\item The jump cancel is the quoter's protection; if it waited for tokens, a budget breach would turn into a pick-off.
\item By the references in its acknowledgements, which arrive before the public feed; if the feed came first, the quoter would take its own order for
someone else's and chase it.
\item Book 7's tape flow replayed as real orders; a jump of the efficient price every 20 seconds on average; an arbitrageur who hears of it first
and takes the stale best level; the quoter hears 20 microseconds later and cancels its stale side.
\item Lit, inverted, bumped and FIFO futures: 14,567 units, 42\% picked off, $+0.03$ ticks, 591 messages; batch 2,567, 0\%, $-1.12$, 330; pro rata
10,884, 47\%, $-0.17$, 474; crypto 7,267, 56\%, $-0.07$, 2,857; event 8,367, 6\%, $-0.52$, 799.
\item Pro rata (size 500: 40,633 units, 58\% picked off, $-0.14$, 419 messages); crypto (16-s pacing: 8,333 units, 191 messages, no rejects);
event (3-tick threshold: 7,900 units, 238 messages, $-0.57$).
\item Up to eight seconds the venue accepts the whole budget, 151 to 154 orders; volume rises as pacing spreads them, peaks at eight seconds and falls
beyond sixteen, where the budget is no longer used.
\item The table of the chapter's result: volume, share picked off, edge and messages on each venue, home and own.
\item A parameter: pro rata, crypto, event. Nothing: lit, inverted (the fee line only), bumped, FIFO futures. Different code: batch and event.
\item The delay applies to the cancel as to the arbitrageur's order; an asymmetric delay that spares cancels would.
\item Price new orders for where the book will be when they land, or send none while the price moves; keep resting orders, which the delay
protects.
\item Quote around a fair value rather than the touch, and size by the batch's imbalance.
\item Write the venue half from the rulebook and specifications; validate with the home quoting half; run the simulator with home and proposed
parameters; the venue's test environment and conformance tests; small size under risk controls; full size once fills, rejects and fees match the
profile.
\item Version it; each venue notice is a new version, validated and replayed before its effective date; the desk head and the venue's owner sign off.
\item Learning a venue's rules before the first order rather than from the losses.
\end{enumerate}
\end{solution}

\begin{solution}{iq:hf:a-venue-playbook:1}
Matching rule, tick and lot, fees and unit, feed, order types, delays, budgets, session and quirks, and the desk's parameters. Before the first
order: the lot (a size the venue refuses), the budget and what a breach does, and the fee sign.
\iqlookfor{rules tied to parameters, not a list of facts.}
\end{solution}

\begin{solution}{iq:hf:a-venue-playbook:2}
The size: allocation follows displayed size, so the home size gets a small share. The risk is that a larger size is also more of what is taken when
the price jumps (58\% against 47\% of volume in the chapter's runs), and that everyone over-quotes in turn.
\iqlookfor{size as the pro-rata parameter, and its cost.}
\end{solution}

\begin{solution}{iq:hf:a-venue-playbook:3}
The order budget: requotes at every book update burn the tokens in bursts, each rejected order is sent again, and the tokens are gone when the market
moves. Pace requotes within the budget, keep cancels outside it, and spend the burst on moves, not on noise.
\iqlookfor{reading rejects as a budget, not a network fault.}
\end{solution}

\begin{solution}{iq:hf:a-venue-playbook:4}
Resting quotes can be cancelled before the delayed taker arrives, so the maker's adverse selection from fast takers falls and it can quote tighter
or bigger; the fast takers lose, and so, in part, do slow takers who wait longer. A maker whose own new orders are delayed pays in stale placements.
\iqlookfor{who pays for the protection.}
\end{solution}

\begin{solution}{iq:hf:a-venue-playbook:5}
At most $b$ orders a second on average, so a requote of two orders every $2/b$ seconds never breaches. A faster pacing can trade more because the
burst $B$ absorbs short spurts and the quotes stay nearer the touch, at the cost of occasional rejects.
\iqlookfor{average rate against burst.}
\end{solution}

\begin{solution}{iq:hf:a-venue-playbook:6}
The maker's bid at the touch executes only if the clearing price is at or below it, which happens when sell orders in the batch exceed buy orders at
that price: its fills are conditioned on selling pressure, so they are followed by lower prices. A rule that quotes around a fair value adjusted for
the batch's expected imbalance, and sizes down when the imbalance is informative, is not adversely selected in the same way.
\iqlookfor{conditioning on the clearing event.}
\end{solution}
